\mychapter{III. METHODOLOGY}{III. METHODOLOGY}

\section{System Topology}

\hs This chapter describes how the system was actually built, from the raw law
documents at one end to the answer shown on the user's screen at the other. The
AI-Powered Legal Assistant for Cambodian Law is organised as an Agentic
Retrieval-Augmented Generation (RAG) pipeline. Rather than sending a question straight
to a language model and hoping for the best, the system passes each query through a
series of decision-making stages, and only the last of these produces the final Khmer
legal answer. Each stage has a narrow, well-defined job, which makes the system easier
to reason about and easier to fix when something goes wrong.

\hs Looking at the topology as a whole helps before examining the parts in isolation. A
user question enters the pipeline and is first classified by intent, so that the system
knows what kind of request it is dealing with. If the question calls for legal
information, it is rewritten into more formal language, used to search the law through a
hybrid retrieval step, and then checked to confirm that the retrieved articles are
actually relevant. Only after that check does the system generate an answer, which is
returned to the frontend together with citations and the ability to open the source law
at the exact page. The figure below shows this flow.

%\fg[0.95\textwidth]{System Architecture Diagram}{fig:system_architecture}
\fg{0.8\textwidth}{System Architecture Diagram}{images/img-ch3/agentic_legal_bot_system_architecture.drawio.png}

\section{Steps of Implementation}

\subsection{Data Collection}

\hs Every retrieval system is only as good as the documents behind it, so the starting
point was the law itself. The knowledge base is built from four official Cambodian
labour laws, combined into a single PDF of 204 pages. These are the Labor Law (393
articles, pages 7--99), the Trade Union Law (100 articles, pages 107--139), the Social
Security Law (100 articles, pages 147--175), and the Minimum Wage Law (29 articles,
pages 191--198), giving 622 articles in total. The documents were taken from official
Cambodian government and legal sources, so that the system answers from the authoritative
text rather than from a paraphrase or summary.

\hs Turning a long PDF into something a retrieval system can use required parsing it
into individual articles. A custom Python script handled this. It located the boundaries
between articles using Khmer article markers such as \textit{មាត្រា}, along with chapter
headings, section headings, and page-number annotations in the text. Each article was
then stored as a structured chunk carrying the metadata fields \texttt{chunk\_id},
\texttt{law\_name}, \texttt{law\_name\_en}, \texttt{article\_number},
\texttt{article\_title}, \texttt{chapter}, \texttt{section}, \texttt{page\_number}, and
\texttt{pdf\_filename}. Keeping this metadata attached to every chunk is what later makes
it possible to show a citation and jump to the right page of the original document.

\begin{table}[H]
\centering
\caption{Laws Indexed in the Knowledge Base}
\label{tab:laws_indexed}
\begin{tabular}{p{5cm} p{4cm} r r}
\toprule
\textbf{Law (Khmer)} & \textbf{Law (English)} & \textbf{Articles} & \textbf{Pages} \\
\midrule
{\khmerfont ច្បាប់ស្តីពីការងារ} & Labor Law & 393 & 7--99 \\
{\khmerfont ច្បាប់ស្តីពីសហជីព} & Trade Union Law & 100 & 107--139 \\
{\khmerfont ច្បាប់ស្តីពីរបបសន្តិសុខសង្គម} & Social Security Law & 100 & 147--175 \\
{\khmerfont ច្បាប់ស្តីពីប្រាក់ឈ្នួលអប្បបរមា} & Minimum Wage Law & 29 & 191--198 \\
\midrule
\textbf{Total} & & \textbf{622} & \textbf{204} \\
\bottomrule
\end{tabular}
\end{table}

\subsection{Data Preprocessing}

\hs Because the system uses two different kinds of search, each article chunk was
prepared in two ways. The first is the raw Khmer article text, which feeds the BM25
keyword index. This text is tokenised with a Khmer Unicode regular-expression tokenizer
that extracts sequences of Khmer characters in the range U+1780--U+17FF. Khmer does not
separate words with spaces, so a tokenizer that understands the script is necessary;
splitting on spaces alone would not work. The second representation is an
\texttt{embed\_text} field, which joins the law name, article number, article title,
and article body into a single string. Giving the embedding model this fuller context,
rather than the bare article body, helps it produce a vector that captures what the
article is actually about.

\hs All 622 chunks were written to a single JSONL file, \texttt{law\_chunks.jsonl},
which serves as the one source of truth for both indexes. Building both the vector index
and the keyword index from the same file avoids the risk of the two drifting out of
sync. The page numbers stored here are the printed page numbers as they appear in the
physical document, which is what allows the frontend to open the PDF at the correct page
later on.

\subsection{Index Construction}

\hs From the preprocessed chunks, two complementary indexes were built, one for meaning
and one for exact words.

\textbf{Dense Vector Index (FAISS):} Each chunk's \texttt{embed\_text} was embedded
using the \texttt{gemini-embedding-001} model with
\texttt{task\_type=RETRIEVAL\_DOCUMENT} and \texttt{output\_dimensionality=768}. The
choice of task type matters here: the embedding model is told that these vectors
represent documents to be retrieved, which aligns them with the query-side embeddings
created later. All 622 vectors were L2-normalised and added to a FAISS
\texttt{IndexFlatIP} index. Because the vectors are normalised, the inner product
computed by this index is equivalent to exact cosine similarity. The index and its chunk
metadata were saved to disk as \texttt{faiss.index} and \texttt{chunks.pkl}.

\textbf{Sparse Keyword Index (BM25):} In parallel, the BM25Okapi algorithm was applied
to the tokenised Khmer corpus. The resulting index, together with the full tokenised
corpus, was stored in \texttt{bm25.pkl} so that it can be loaded directly at runtime
without rebuilding. Keeping a sparse index alongside the dense one is a deliberate
decision, motivated by the low-resource nature of Khmer discussed earlier: when the
embeddings miss a rare legal term, exact keyword matching can still find it.

\begin{table}[H]
\centering
\caption{Index Configuration}
\label{tab:index_config}
\begin{tabular}{lll}
\toprule
\textbf{Index} & \textbf{Method} & \textbf{Configuration} \\
\midrule
Dense & FAISS IndexFlatIP & 622 vectors, dim=768, cosine similarity \\
Sparse & BM25Okapi & 622 documents, Khmer regex tokeniser \\
\bottomrule
\end{tabular}
\end{table}

\subsection{Agent Pipeline Design}

\hs The heart of the system is a LangGraph \texttt{StateGraph} that connects five
processing nodes in sequence. The design follows a simple rule: every node receives the
full agent state as a dictionary and returns an updated copy of that state. Nothing is
hidden in global variables, so the information available to each node is always explicit
and easy to trace. The graph is compiled once when the server starts and then reused as
a singleton for every request, which avoids the cost of rebuilding it each time and keeps
behaviour consistent across users. The sections that follow describe what each node does
and why it is there.

\addcontentsline{toc}{subsection}{Intent Classifier}
\subsubsection*{Intent Classifier}

\hs Not every message a user sends is a legal question, and treating all of them the same
way would waste effort and sometimes give odd answers. The intent classifier node, built
on \texttt{gemini-2.5-flash}, sorts each query into one of five categories:
\texttt{legal\_qa}, \texttt{article\_lookup}, \texttt{definition},
\texttt{recent\_news}, and \texttt{greeting}. To keep this step fast and predictable, the
model is configured with \texttt{response\_mime\_type=application/json},
\texttt{temperature=0.0}, and \texttt{thinking\_budget=0}, so it returns a clean
structured result without spending time on internal reasoning. The classified intent then
decides the route: \texttt{legal\_qa} and \texttt{definition} go through the full RAG
pipeline, \texttt{article\_lookup} is sent to a direct article lookup that matches by
number, and \texttt{greeting} or \texttt{recent\_news} skip retrieval entirely and go
straight to a simple response. Routing this way means the expensive retrieval machinery
only runs when it is actually needed.

\addcontentsline{toc}{subsection}{Query Rewriter}
\subsubsection*{Query Rewriter}

\hs The query rewriter addresses the central difficulty identified in the background: the
words a person uses are rarely the words the law uses. Running on
\texttt{gemini-2.5-flash}, this node takes an informal Khmer question and expands it into
formal legal vocabulary. Its system prompt asks the model to work out which of the four
laws the question relates to, to add the formal Khmer legal terms that the relevant
articles are likely to contain, and to return a rewritten query better suited to
retrieval. In practice this step often makes the difference between finding the right
article and missing it, because it gives both the dense and the sparse search a better
target to match against.

\addcontentsline{toc}{subsection}{Hybrid Retrieval with RRF}
\subsubsection*{Hybrid Retrieval with RRF}

\hs The retrieval node is where the two indexes come together. The rewritten query is
embedded with \texttt{gemini-embedding-001} using
\texttt{task\_type=RETRIEVAL\_QUERY}, the query-side counterpart to the document task
type used during indexing, and then L2-normalised before FAISS returns its top-20
candidates by cosine similarity. At the same time, the query is tokenised with the same
Khmer regex tokeniser and scored against all 622 BM25 entries, producing another top-20
list. The system now has two ranked lists that may not agree, and the question is how to
combine them fairly. Reciprocal Rank Fusion solves this without needing to compare the
two systems' raw scores, which are not on the same scale. With a smoothing constant
$k = 60$, each document's fused score is:

\begin{equation}
\text{RRF}(d) = \sum_{r \in R} \frac{1}{k + \text{rank}_r(d)}
\end{equation}

\hs A document that ranks highly in either list, or moderately in both, rises to the
top. The five highest-scoring chunks are then taken forward as the retrieved context for
the grounding check and answer generation. Five was chosen as a balance: enough context
to answer most questions, but not so much that the language model is overwhelmed with
loosely related articles.

\addcontentsline{toc}{subsection}{Grounding Verification}
\subsubsection*{Grounding Verification}

\hs Retrieval always returns something, even when nothing relevant exists, so the
grounding node acts as a gatekeeper before any answer is written. It passes the first 500
characters of each retrieved chunk to \texttt{gemini-2.5-flash} and asks for a simple
binary judgement, returned as \texttt{\{"is\_relevant": true/false\}}. If the answer is
negative and fewer than two retries have been used, the pipeline loops back to the query
rewriter, which produces a fresh query for another retrieval attempt. If the retry limit
is reached, the system proceeds to generation with whatever it has, rather than looping
forever. This retry-then-proceed design borrows the spirit of self-reflection from the
literature while keeping the behaviour bounded and predictable.

\addcontentsline{toc}{subsection}{Answer Generation}
\subsubsection*{Answer Generation}

\hs The final stage produces the answer, and it does so in two separate language-model
calls rather than one. The reason is practical: asking a single call to write fluent
Khmer prose and, at the same time, emit perfectly structured citation data tends to
degrade both. Splitting the work lets each call do one thing well.

\textbf{Step 1 — Khmer Answer:} The retrieved article context is combined with the user
query and sent to \texttt{gemini-2.5-flash}, with a system prompt instructing it to
answer in formal Khmer, to cite articles using inline \texttt{[N]} markers, and to keep
the response to roughly 200 words covering two or three key points. The aim is an answer
that is accurate and readable rather than exhaustive.

\textbf{Step 2 — Citation Extraction:} A second, lightweight call reads the completed
answer and extracts the citations as a JSON array, where each object carries the
\texttt{article\_number}, \texttt{law\_name}, \texttt{page\_number}, and
\texttt{ref\_num}.

\textbf{Step 3 — Citation Filter:} Finally, a small post-processing step keeps only those
citations whose reference number \texttt{[N]} actually appears in the answer text. Any
article that was retrieved but not used is dropped, so the citation chips shown to the
user reflect only the sources the answer genuinely relied on.

%\imgplaceholder[0.95\textwidth]{Agent Pipeline Flow Diagram}{fig:pipeline_flow}
\fg{0.8\textwidth}{Agent Pipeline Flow Diagram}{images/img-ch3/system_architecture_layer.drawio.png}

\subsection{System Integration}

\hs With the agent pipeline in place, the remaining work was to wrap it in a usable
application. The backend is a FastAPI service run by Uvicorn. Its main endpoint,
\texttt{POST /api/chat}, receives the user's message and session ID, runs the full agent
graph inside a thread executor so that the synchronous pipeline does not block the async
event loop, and then streams the finished answer back to the browser using Server-Sent
Events (SSE). The stream follows a fixed sequence of events: status messages while the
system works, the answer delivered in small text chunks, a citations event, and finally
a done event that tells the client the response is complete.

\hs The frontend is a single-page application built with React 19 and Vite. It reads the
SSE stream through the Fetch API's \texttt{ReadableStream}, showing the Khmer answer as
it arrives so the user is not left staring at a blank screen. Once streaming finishes, the
citation chips appear, and clicking one opens a PDF viewer panel that renders the source
law at the exact article page using pdfjs-dist. This direct link from answer to source is
what lets a user verify the system's claims for themselves, which matters a great deal in
a legal setting.

\hs Conversation history is stored in a Neon serverless PostgreSQL database. Each session
keeps its user messages and assistant responses, with the citation data saved as JSON, in
the \texttt{chat\_messages} table. This makes it possible to reload a past conversation
later. The full set of endpoints is summarised in Table~\ref{tab:api_endpoints}.

\begin{table}[H]
\centering
\caption{API Endpoints}
\label{tab:api_endpoints}
\begin{tabular}{llp{7cm}}
\toprule
\textbf{Method} & \textbf{Endpoint} & \textbf{Description} \\
\midrule
POST & \texttt{/api/chat} & Main chat endpoint — runs agent pipeline, returns SSE stream \\
POST & \texttt{/api/chat/session} & Creates a new chat session, returns session ID \\
GET  & \texttt{/api/chat/\{id\}/history} & Retrieves message history for a session \\
GET  & \texttt{/api/documents/\{filename\}} & Serves PDF files with whitelist protection \\
GET  & \texttt{/api/health} & Health check endpoint \\
\bottomrule
\end{tabular}
\end{table}

\subsection{Evaluation Metrics}

\hs Measuring whether the system works required metrics suited to a document-retrieval
question answering task, rather than the accuracy figures used for classification. Let
$R_k$ be the set of top-$k$ retrieved articles for a query, and $A$ the set of articles
that genuinely answer it.

\textbf{Retrieval Precision@K} asks a basic question: did the search find a correct
article at all? It measures the proportion of queries for which at least one expected
article appears in the top-$K$ results:

\begin{equation}
\text{Precision@K} = \frac{1}{|Q|} \sum_{q \in Q} \mathbf{1}[A_q \cap R_k(q) \neq \emptyset]
\end{equation}

\textbf{Citation Accuracy} goes one step further and checks the final answer rather than
the retrieval. It measures the proportion of queries for which a correct article appears
as a cited reference in the answer the user actually sees:

\begin{equation}
\text{Citation Accuracy} = \frac{1}{|Q|} \sum_{q \in Q} \mathbf{1}[A_q \cap C(q) \neq \emptyset]
\end{equation}

\noindent where $C(q)$ is the set of articles cited inline in the generated answer. The
gap between this metric and Precision@K reveals cases where the right article was found
but not used.

\textbf{Grounding Rate} reports how often the grounding node judged the retrieved
articles sufficient, which gives a sense of retrieval quality from the system's own point
of view:

\begin{equation}
\text{Grounding Rate} = \frac{|\{q \in Q : \text{grounded}(q) = \text{true}\}|}{|Q|}
\end{equation}

\subsection{Experimental Setup}

\hs The system was developed and tested on a MacBook Air running macOS with an Apple
Silicon processor. The backend runs in a Python 3.11 virtual environment. All language
model and embedding calls go through Google Vertex AI, using \texttt{gemini-2.5-flash}
for generation and classification and \texttt{gemini-embedding-001} for both document
and query embeddings. Authentication to Vertex AI uses Application Default Credentials
through \texttt{gcloud auth application-default login}, which avoids storing an API key in
the code. The database is hosted on Neon's serverless PostgreSQL and accessed with
\texttt{asyncpg}. The full environment is listed in Table~\ref{tab:exp_setup}.

\begin{table}[H]
\centering
\caption{Experimental Environment}
\label{tab:exp_setup}
\begin{tabular}{ll}
\toprule
\textbf{Component} & \textbf{Specification} \\
\midrule
Operating System & macOS (Apple Silicon) \\
Backend Language & Python 3.11 \\
Frontend Language & JavaScript (React 19, Vite 8) \\
LLM API & Google Vertex AI \\
Generation Model & gemini-2.5-flash \\
Embedding Model & gemini-embedding-001 (dim=768) \\
Vector Index & FAISS IndexFlatIP (622 vectors) \\
Keyword Index & BM25Okapi (622 documents) \\
Database & Neon PostgreSQL (serverless) \\
Evaluation Dataset & 30 questions, 10 categories \\
\bottomrule
\end{tabular}
\end{table}

\subsection{Deployment}

\hs Beyond local development, the system was prepared for deployment so that it could run
as a real service. The FastAPI backend is served by Uvicorn and can be hosted on any
suitable cloud platform, while the React frontend is built with Vite into static files
that can be served on their own. During development the two are connected through Vite's
proxy, which forwards every \texttt{/api/*} request to the FastAPI server on port 8000.
CORS is configured to accept requests only from the frontend origin. The combined law PDF
is served through a whitelisted file endpoint, which prevents path-traversal attacks by
refusing to serve any file outside the approved list. Because all session state lives in
the PostgreSQL database rather than in the server's memory, the backend remains stateless
between requests, which keeps it simple to restart or scale.

%\imgplaceholder[0.85\textwidth]{System Deployment Diagram}{fig:deployment}
\fg{0.8\textwidth}{System Deployment Diagram}{images/img-ch3/System-Deployment-Diagram.drawio.png}
% ─────────────────────────────────────────────────────────────────────
\newpage
\section{Tools}

\hs This section lists the main software tools used during development. Each played a
specific role in building, testing, or documenting the system.

\subsection{Visual Studio Code}
\textbf{Visual Studio Code (VS Code)} is a free, open-source code editor developed by
Microsoft. It supports syntax highlighting, intelligent code completion, debugging, and
version control integration. VS Code was the primary development environment for writing
and testing both the Python backend and the React frontend of this project.
%\imgplaceholder[0.23\textwidth]{Visual Studio Code Logo}{fig:vscode_logo}
\fg{0.2\textwidth}{Visual Studio Code Logo}{images/img-ch3/Visual-Studio-Code.png}

\subsection{Draw.io}
\textbf{Draw.io} is a free web-based diagramming application for creating flowcharts,
architecture diagrams, and system topology figures. It was used to design the agent
pipeline flow diagram and the system architecture diagram presented in this thesis.
%\imgplaceholder[0.23\textwidth]{Draw.io Logo}{fig:drawio_logo}
\fg{0.2\textwidth}{Draw.io Logo}{images/img-ch3/drawio.png}

\subsection{Postman}
\textbf{Postman} is an API testing platform that lets developers send HTTP requests and
inspect the responses. It was used to test and validate all FastAPI endpoints during
backend development, including the chat, session, history, and document endpoints, before
they were connected to the frontend.
%\imgplaceholder[0.23\textwidth]{Postman Logo}{fig:postman_logo}
\fg{0.3\textwidth}{Postman Logo}{images/img-ch3/Postman.png}

\subsection{Google Cloud Console}
\textbf{Google Cloud Console} is the web-based management interface for Google Cloud
Platform. It was used to manage the Vertex AI project, configure Application Default
Credentials, enable the Vertex AI API, and monitor API usage and billing for the
\texttt{gemini-2.5-flash} and \texttt{gemini-embedding-001} models.
%\imgplaceholder[0.23\textwidth]{Google Cloud Console Logo}{fig:gcp_logo}
\fg{0.3\textwidth}{Google Cloud Console Logo}{images/img-ch3/Google-Cloud-Logo.png}

\subsection{Overleaf}
\textbf{Overleaf} is a collaborative online \LaTeX{} editor for writing and compiling
academic documents. It was used to write and format this thesis, providing real-time PDF
compilation, version history, and multi-file project management.
%\imgplaceholder[0.23\textwidth]{Overleaf Logo}{fig:overleaf_logo}
\fg{0.2\textwidth}{Overleaf Logo}{images/img-ch3/Overleaf-Logo.png}
% ─────────────────────────────────────────────────────────────────────
\section{Programming Languages}

\hs Two programming languages were used, one for each side of the application.

\subsection{Python}
\textbf{Python} is a high-level, interpreted language known for its simplicity and
versatility. It is the primary language of the backend, used to implement the FastAPI
server, the LangGraph agent pipeline, FAISS retrieval, BM25 indexing, and all data
processing scripts. Python 3.11 was used throughout the project.
%\imgplaceholder[0.23\textwidth]{Python Logo}{fig:python_logo}
\fg{0.2\textwidth}{Python Logo}{images/img-ch3/python.jpg}

\subsection{JavaScript}
\textbf{JavaScript} is a high-level scripting language used mainly for web development.
Here it powers the React frontend, including the chat interface, the SSE stream consumer,
the citation chip components, and the PDF viewer integration with pdfjs-dist.
%\imgplaceholder[0.23\textwidth]{JavaScript Logo}{fig:js_logo}
\fg{0.2\textwidth}{JavaScript Logo}{images/img-ch3/Javascript-Logo.png}
% ─────────────────────────────────────────────────────────────────────

\section{Libraries and Frameworks}

\hs The system relies on a number of open-source libraries and frameworks, each handling
a specific part of the work.

\subsection{FastAPI}
\textbf{FastAPI} is a modern, high-performance Python web framework for building APIs with
automatic OpenAPI documentation. Its support for asynchronous request handling through
Python's \texttt{asyncio} makes it well suited to serving streaming SSE responses and to
handling several user sessions at once. FastAPI was used to implement all backend API
endpoints in this project.
%\imgplaceholder[0.23\textwidth]{FastAPI Logo}{fig:fastapi_logo}
\fg{0.3\textwidth}{FastAPI Logo}{images/img-ch3/FastAPI-Logo.png}

\subsection{LangGraph}
\textbf{LangGraph} is a library for building stateful, multi-actor applications with
language models using directed graphs. It extends LangChain with support for cyclic
graphs and conditional edges, which is exactly what the retry loop and intent-based
routing in this system require. The \texttt{StateGraph} class was used to define and
compile the five-node agent pipeline as a singleton reused across all requests.
%\imgplaceholder[0.23\textwidth]{LangGraph Logo}{fig:langgraph_logo}
\fg{0.3\textwidth}{LangGraph Logo}{images/img-ch3/Langgraph-seeklogo.png}

\newpage
\subsection{FAISS}
\textbf{FAISS} (Facebook AI Similarity Search) is a library developed by Meta AI for
efficient similarity search over large collections of dense vectors. It supports both
exact and approximate nearest-neighbour search on CPU and GPU. In this project,
\texttt{faiss-cpu} with \texttt{IndexFlatIP} provides exact cosine similarity search over
the 622 article embeddings, each of dimension 768.

\subsection{rank-bm25}
\textbf{rank-bm25} is a Python implementation of the BM25Okapi ranking algorithm for
sparse, keyword-based document retrieval. It was used to build and query the BM25 keyword
index over the 622 Khmer legal article chunks, providing the exact-term matching that
complements FAISS dense retrieval.

\subsection{Google Generative AI SDK (google-genai)}
\textbf{google-genai} (v2.4.0) is the official Python SDK for accessing Google's
generative models, including Gemini. It provides text generation, structured JSON output,
and document embedding through a single client interface. Every language model call in
this project — intent classification, query rewriting, grounding checking, answer
generation, and citation extraction — uses this SDK through the Vertex AI backend.
%\imgplaceholder[0.23\textwidth]{Google Gemini Logo}{fig:gemini_logo}
\fg{0.3\textwidth}{Google Gemini Logo}{images/img-ch3/Google-Gemini-logo.png}

\subsection{React}
\textbf{React} (v19) is a JavaScript library for building user interfaces from
composable components. It was used to implement the frontend of the legal chatbot,
including the chat window, message bubbles, citation chips, typing indicator, welcome
screen, and collapsible sidebar.
%\imgplaceholder[0.23\textwidth]{React Logo}{fig:react_logo}
\fg{0.2\textwidth}{React Logo}{images/img-ch3/ReactJs-logo.png}

\newpage
\subsection{pdfjs-dist}
\textbf{pdfjs-dist} (v5.7) is the Mozilla PDF.js library distributed as an npm package.
It exposes a JavaScript API for rendering PDF documents onto an HTML canvas. In this
project it powers the in-browser PDF viewer that opens the source law at the exact page of
a cited article whenever a citation chip is clicked.
%\imgplaceholder[0.23\textwidth]{PDF.js Logo}{fig:pdfjs_logo}
\fg{0.2\textwidth}{PDF.js Logo}{images/img-ch3/PdfJs-Logo.png}

\subsection{asyncpg}
\textbf{asyncpg} is a high-performance asynchronous PostgreSQL client for Python. It
offers a raw SQL interface without an ORM, using positional placeholders such as
\texttt{\$1} and \texttt{\$2} for parameterised queries. It handles all database
operations in this project, including session creation, message storage, and history
retrieval against the Neon serverless PostgreSQL database.

\subsection{Pydantic}
\textbf{Pydantic} is a Python data-validation library for defining and validating data
models and settings. Through \texttt{pydantic-settings}, it manages the application's
configuration — database URLs, model names, retrieval parameters, and keys — loading them
from environment variables into a single typed \texttt{Settings} object.

\subsection{Vite}
\textbf{Vite} (v8) is a modern frontend build tool offering fast hot module replacement
during development and optimised production builds. It was used to scaffold and build the
React frontend, and its built-in proxy was configured to forward all \texttt{/api/*}
requests to the FastAPI backend during development.
%\imgplaceholder[0.23\textwidth]{Vite Logo}{fig:vite_logo}
\fg{0.2\textwidth}{Vite Logo}{images/img-ch3/Vite-Logo.jpeg}

% ─────────────────────────────────────────────────────────────────────
\newpage
\section{Infrastructure, Deployment, and Services}

\hs Beyond the core libraries, the system depends on a number of infrastructure tools
and external services to run as a real, publicly accessible application. These handle
hosting, containerisation, secure access, authentication, and data storage.

\subsection{Docker}
\textbf{Docker} is a platform for packaging an application and its dependencies into
portable containers that run the same way on any machine. In this project, the backend
and frontend were each given a Dockerfile, and a \texttt{docker-compose.yml} file ties
them together so that the whole system can be started with a single command. Containerising
the application made the move from local development to the production server reliable,
since the server runs exactly the same images that were tested locally.
%\imgplaceholder[0.23\textwidth]{Docker Logo}{fig:docker_logo}
\fg{0.3\textwidth}{Docker Logo}{images/img-ch3/docker-logo-ocean-blue.png}

\subsection{Amazon Web Services EC2}
\textbf{Amazon EC2} (Elastic Compute Cloud) provides resizable virtual servers in the
cloud. The production system is hosted on a \texttt{t3.small} EC2 instance running Ubuntu,
with a fixed Elastic IP address so that the public address does not change when the server
restarts. The Docker containers for the backend and frontend run on this instance, making
the application reachable from anywhere over the internet.
%\imgplaceholder[0.23\textwidth]{AWS EC2 Logo}{fig:ec2_logo}
\fg{0.1\textwidth}{AWS EC2 Logo}{images/img-ch3/EC2.png}

\subsection{Nginx}
\textbf{Nginx} is a high-performance web server and reverse proxy. In the production
deployment it serves the static React build, forwards all \texttt{/api/*} requests to the
FastAPI backend container, and terminates HTTPS connections. Its configuration was tuned
to support Server-Sent Events by disabling response buffering, so that answers still
stream to the user token by token in production.
%\imgplaceholder[0.23\textwidth]{Nginx Logo}{fig:nginx_logo}
\fg{0.1\textwidth}{Nginx Logo}{images/img-ch3/nginx.png}

\subsection{Let's Encrypt and Certbot}
\textbf{Let's Encrypt} is a free, automated certificate authority, and \textbf{Certbot}
is the client used to obtain and install its certificates. Together they provide the
HTTPS certificate for the production domain, so that traffic between users and the server
is encrypted. The certificate was issued and configured through a one-time setup script
that runs Certbot in a container alongside Nginx.
%\imgplaceholder[0.23\textwidth]{Let's Encrypt Logo}{fig:letsencrypt_logo}
\fg{0.2\textwidth}{Let's Encrypt Logo}{images/img-ch3/cerbot-logo.png}

\subsection{Clerk}
\textbf{Clerk} is a managed authentication and user-management service. It handles user
sign-up and sign-in, including Google sign-in, and issues signed session tokens. In this
project, the frontend uses Clerk to manage the sign-in flow, while the backend verifies
the Clerk session token on each request to identify the user. Authentication is optional:
anonymous visitors may ask a limited number of questions, after which they are prompted to
sign in to continue, and signed-in users gain access to their saved conversation history.
%\imgplaceholder[0.23\textwidth]{Clerk Logo}{fig:clerk_logo}
\fg{0.2\textwidth}{Clerk Logo}{images/img-ch3/clerk-logo.png}

\subsection{Neon (Serverless PostgreSQL)}
\textbf{Neon} is a serverless PostgreSQL platform. It hosts the project's database, which
stores chat sessions, messages, and citation data, as well as the per-user association
that powers conversation history. Being serverless, it requires no database server to be
managed manually and scales automatically with demand.
%\imgplaceholder[0.23\textwidth]{Neon Logo}{fig:neon_logo}
\fg{0.2\textwidth}{Neon Logo}{images/img-ch3/neon-logo.png}

\subsection{Git and GitHub}
\textbf{Git} is a distributed version-control system, and \textbf{GitHub} is the hosting
platform used for the project's repository. All source code was tracked with Git, and
GitHub served both as a backup and as the means of deploying to the production server,
which pulls the latest code directly from the repository.
%\imgplaceholder[0.23\textwidth]{GitHub Logo}{fig:github_logo}
\fg{0.1\textwidth}{GitHub Logo}{images/img-ch3/github.png}

